% Draft only. Not included by conference_101719_1.tex.
% Branch analysis/failure-modes, cut from overleaf/main at bbd5bd8.
% Reproducing commands for every number are in
% docs/failure_modes_and_threshold_sensitivity.md.

\subsection{Failure modes and threshold sensitivity}

\subsubsection{Mode classification}
The flood-vehicle literature distinguishes instability modes by force balance rather
than by any kinematic cutoff. AR\&R Project 10 identifies two recognised hydrodynamic
mechanisms, floating and frictional sliding, and explicitly excludes toppling from
its analysis on the grounds that it is restricted to vehicles already sliding or
floating over uneven terrain. Shah \emph{et al.} likewise define floating by buoyancy
and lift exceeding vehicle weight. Because no source supplies a displacement or
rotation threshold, the cutoffs used here are \emph{operational} detection tolerances
for classifying a simulation trace, not physical instability criteria, and each is
swept rather than fixed.

Applying a drift cutoff of $0.05$\,m, a roll cutoff of $1.0^\circ$, and a vertical-rise
cutoff of $0.01$\,m to the 17-run sweep yields 15 SLIDE and 2 TOPPLE events, with no
STUCK and no FLOAT. Sliding is the dominant mode throughout. The two toppling events
occur at the extremes of the sweep, the deepest condition ($0.45$\,m) and the fastest
($3.0$\,m/s), consistent with AR\&R's characterisation of toppling as a consequence of
prior instability rather than an independent mode. The TOPPLE count is stable at two
across roll cutoffs spanning $0.51^\circ$ to $2.43^\circ$, because the measured roll
magnitudes separate into two well-defined groups with an empty interval between them.

Two null results require care. STUCK is empty because the smallest drift observed,
$0.0578$\,m, already exceeds the $0.05$\,m cutoff; every run moves. FLOAT is empty at
every rise cutoff tested, but this is not evidence against flotation. The minimum
vertical coordinate of the vehicle equals three grid cells in all 17 runs, to within
$4\times10^{-8}$\,m, which is the solver's inward domain padding. The vertical channel
is saturated at a boundary, so flotation is unobservable in this configuration rather
than absent from it.

\subsubsection{Threshold sensitivity}
The $0.05$\,m drift threshold has no empirical footing in the literature, so its
influence is characterised by sweeping it. Runs are matched to the L1 scenario grid by
exact agreement in depth and velocity, and each run is compared against the L1 verdict
for its own vehicle class. Fourteen of 17 runs match; the three excluded lie at depths
of $0.25$, $0.35$, and $0.45$\,m, which do not appear on the $0.1$\,m scenario grid.
The matched runs occupy six distinct scenario cells, all at $0.30$\,m depth, so the
comparison characterises one depth column rather than the full hazard plane.

Pooled over the matched runs, no threshold reconciles the two levels: agreement peaks
at $92.9\%$ for thresholds in $[0.26, 0.35]$\,m and never reaches unity. At the
$0.05$\,m operating point agreement is $42.9\%$, and the curve is flat below
$0.0578$\,m and varies by only $7.1$ percentage points across $[0.03, 0.07]$\,m, so the
verdict at that operating point is insensitive to the threshold. That insensitivity
reflects the threshold sitting an order of magnitude below the data rather than any
robustness of the comparison.

This pooled result does not survive stratification, and the stratified result is the
substantive one. Separating the matched runs by grid resolution, a threshold achieving
complete agreement exists at every resolution: $[0.257, 0.351)$\,m at $n_g=48$,
$[0.314, 0.659)$\,m at $n_g=64$, and $[0.156, 0.269)$\,m at $n_g=96$. The three
intervals are individually non-empty and pairwise disjoint. Within any single
resolution, simulated drift orders the conditions exactly as the L1 hazard product
does; the pooled disagreement arises entirely from resolution-dependent scatter. At
fixed physical conditions ($D=0.30$\,m, $V=1.5$\,m/s) drift varies by a factor of
$2.0$ to $2.5$ across the three resolutions, and non-monotonically for the lightest
vehicle. We therefore report no structural divergence between L1 and L2 from this
dataset. The discrepancy is a convergence artifact, and resolving it requires a grid
refinement study rather than a reinterpretation of the abstraction ladder.

\subsubsection{Violation magnitude}
Reporting only whether a criterion is crossed discards how far. Every run in the sweep
exceeds $0.05$\,m, by factors ranging from $1.16$ to $26.8$. No run falls within
$\pm10\%$ of that boundary; the nearest, at $0.5$\,m/s, exceeds it by $15.6\%$. The
binary verdict therefore misrepresents no individual run here, but only because the
threshold lies far below the observed range.

\subsubsection{Temporal convergence}
Per-frame traces show that drift is non-monotonic in 14 of 17 runs, reaching an
interior maximum and then decreasing. Five runs, the two deepest and the three fastest,
have not converged by the final frame and are re-accelerating at truncation, with
terminal speed exceeding twice the minimum speed reached after the drift peak. For
these five the reported final displacement understates the peak excursion by $8.3\%$ to
$24.9\%$, the worst case reaching a peak of $1.78$\,m at frame 56 before returning to
$1.34$\,m at frame 90. The reversal is consistent with wave reflection from the closed
domain. For the five severe conditions the reported displacement is neither a converged
steady state nor a peak, and verdicts derived from it should be treated as provisional
pending longer integration or an open outflow boundary. The remaining 12 runs plateau
well before truncation.

\subsubsection{External consistency}
An independent three-dimensional CFD study of a full-scale medium-size passenger
vehicle reports flotation at $0.38$\,m depth and sliding above a depth-velocity product
of $0.36$\,m$^2$/s, over a validated Froude range of $0.09$ to $2.46$. All nine
distinct conditions in the present sweep fall within that range, spanning $Fr=0.294$ to
$Fr=1.765$, which places this work inside an independently validated envelope.
Substituting $0.36$\,m$^2$/s for the AR\&R small-passenger value of $0.30$\,m$^2$/s
alters 3 of 70 scenario verdicts under a pure product rule, indicating that the ladder's
L1 conclusions are not strongly sensitive to which of the two thresholds is adopted.
